\begin{abstract}

This paper studies local, sensor-driven navigation for unmanned surface vehicles operating in unknown, time-varying current fields with dynamic obstacles. Conventional value-based deep reinforcement learning optimizes expected return and can underestimate rare but catastrophic failures caused by current--obstacle interactions. We propose a Distributional-Risk Planner based on Implicit Quantile Networks (IQN) that models the full return distribution and applies a CVaR-style risk distortion to emphasize tail outcomes during action selection. To balance safety and efficiency across heterogeneous scenes, we further introduce an adaptive risk mechanism that increases risk aversion online as the vehicle approaches sensed obstacles, while remaining near risk-neutral in open water. An energy-related penalty is incorporated into the reward to encourage smoother, energy-efficient control, and we report an energy consumption analysis. Experiments demonstrate that our method dominates the performance trade-off in dense dynamic environments (10 obstacles): Adaptive IQN achieves a \textbf{100\% success rate} (outperforming PPO's 95--98\%) while maintaining superior efficiency (surpassing PPO's speed and energy metrics). Compared to classical APF which can ensure safety but suffers from extreme inefficiency (taking >2x longer), the proposed approach provides a superior balance of safety, speed, and energy.
\end{abstract}

\begin{IEEEkeywords}
Unmanned surface vehicle, implicit quantile network, conditional value-at-risk, dynamic environment
\end{IEEEkeywords}